\documentclass[tikz,border=8pt]{standalone}
\usepackage{tikz}
\usepackage{ctex}
\usetikzlibrary{arrows.meta,positioning,fit,calc,shapes.geometric,backgrounds}
% Shared visual language for LFS architecture diagrams.
\RequirePackage{../../mymath}

\definecolor{LFSBlue}{HTML}{2563EB}
\definecolor{LFSBlueSoft}{HTML}{DBEAFE}
\definecolor{LFSGreen}{HTML}{15803D}
\definecolor{LFSGreenSoft}{HTML}{DCFCE7}
\definecolor{LFSOrange}{HTML}{C2410C}
\definecolor{LFSOrangeSoft}{HTML}{FFEDD5}
\definecolor{LFSRed}{HTML}{B91C1C}
\definecolor{LFSRedSoft}{HTML}{FEE2E2}
\definecolor{LFSPurple}{HTML}{7E22CE}
\definecolor{LFSPurpleSoft}{HTML}{F3E8FF}
\definecolor{LFSGray}{HTML}{475569}
\definecolor{LFSGraySoft}{HTML}{F1F5F9}

\tikzset{
  lfs picture/.style={
    font=\small,
    node distance=8mm and 12mm,
    >=Latex,
    every path/.style={line cap=round, line join=round}
  },
  block/.style={
    draw=LFSBlue,
    fill=LFSBlueSoft,
    rounded corners=2pt,
    line width=0.8pt,
    align=center,
    inner xsep=7pt,
    inner ysep=5pt,
    minimum height=10mm
  },
  compute/.style={
    block,
    draw=LFSPurple,
    fill=LFSPurpleSoft
  },
  storage/.style={
    block,
    draw=LFSGreen,
    fill=LFSGreenSoft
  },
  interface/.style={
    block,
    draw=LFSOrange,
    fill=LFSOrangeSoft
  },
  risk/.style={
    block,
    draw=LFSRed,
    fill=LFSRedSoft
  },
  neutral/.style={
    block,
    draw=LFSGray,
    fill=LFSGraySoft
  },
  decision/.style={
    diamond,
    draw=LFSOrange,
    fill=LFSOrangeSoft,
    line width=0.8pt,
    align=center,
    inner sep=2pt,
    aspect=2
  },
  group/.style={
    draw=LFSGray,
    rounded corners=3pt,
    dashed,
    line width=0.7pt,
    inner sep=6pt
  },
  arrow/.style={
    -{Latex[length=2.4mm,width=1.6mm]},
    draw=LFSGray,
    line width=0.9pt
  },
  data arrow/.style={
    arrow,
    draw=LFSBlue
  },
  control arrow/.style={
    arrow,
    draw=LFSOrange,
    dashed
  },
  residual/.style={
    arrow,
    draw=LFSGreen,
    rounded corners=4pt
  },
  label text/.style={
    font=\scriptsize,
    text=LFSGray,
    fill=white,
    inner sep=1.5pt
  },
  title/.style={
    font=\bfseries,
    text=LFSGray,
    align=center
  }
}

\begin{document}
\begin{tikzpicture}[lfs picture]
  \node[storage] (source) {对象存储\\并行文件系统};
  \node[compute, right=15mm of source] (cpu) {CPU 解析\\Tokenizer};
  \node[compute, right=15mm of cpu] (gpu) {GPU / Accelerator\\训练或推理 Kernel};
  \node[compute, above=14mm of gpu] (ranks) {其他 Rank};
  \node[storage, below=15mm of gpu] (checkpoint) {持久化存储\\Checkpoint};
  \node[neutral, right=17mm of gpu] (telemetry) {监控与日志\\指标 / Trace};
  \node[interface, below=15mm of telemetry] (planner) {容量模型\\吞吐、内存、网络、恢复窗口};

  \draw[data arrow] (source) -- node[label text,above]{数据吞吐} (cpu);
  \draw[data arrow] (cpu) -- node[label text,above]{Batch} (gpu);
  \draw[data arrow, <->] (gpu) -- node[label text,right]{Collective} (ranks);
  \draw[data arrow] (gpu) -- node[label text,right]{Checkpoint} (checkpoint);
  \draw[data arrow, <->] (gpu) -- (telemetry);
  \draw[control arrow] (source) |- (planner);
  \draw[control arrow] (checkpoint) -- (planner);
  \draw[control arrow] (telemetry) -- (planner);

  \begin{scope}[on background layer]
    \node[group, fit=(source)(cpu), label={[title]above:主机与数据通路}] {};
    \node[group, fit=(gpu)(ranks), label={[title]above:加速器与通信域}] {};
    \node[group, fit=(checkpoint)(telemetry)(planner), label={[title]below:持久化与运维域}] {};
  \end{scope}
\end{tikzpicture}
\end{document}
